% Created 2024-07-29 Mon 20:32
% Intended LaTeX compiler: pdflatex
\documentclass[11pt]{article}
\usepackage[utf8]{inputenc}
\usepackage[T1]{fontenc}
\usepackage{graphicx}
\usepackage{longtable}
\usepackage{wrapfig}
\usepackage{rotating}
\usepackage[normalem]{ulem}
\usepackage{amsmath}
\usepackage{amssymb}
\usepackage{capt-of}
\usepackage{hyperref}
\author{Javier Pacheco}
\date{\today}
\title{Utilities}
\hypersetup{
 pdfauthor={Javier Pacheco},
 pdftitle={Utilities},
 pdfkeywords={},
 pdfsubject={Usueful scripts for emacs.},
 pdfcreator={Emacs 29.4 (Org mode 9.7.8)}, 
 pdflang={English}}
\begin{document}

\maketitle
\tableofcontents

\section{Doom Insert item.}
\label{sec:org1449028}
This script wast taken from \emph{doomemacs}, it basicaly insert a item such a heading, subheading, checkbox, etc. based on the item above.

\begin{minted}[linenos=true,numbersep=2pt,breaklines=true,frame=leftline,framerule=2pt,labelposition=bottomline,bgcolor=jpyellow!70]{common-lisp}
;;; javier_pacheco utilities.el --- Some usefull utulities  -*- lexical-binding: t; -*-
(defun +org--insert-item (direction)
  (let ((context (org-element-lineage
                  (org-element-context)
                  '(table table-row headline inlinetask item plain-list)
                  t)))
    (pcase (org-element-type context)
      ;; Add a new list item (carrying over checkboxes if necessary)
      ((or `item `plain-list)
       (let ((orig-point (point)))
         ;; Position determines where org-insert-todo-heading and `org-insert-item'
         ;; insert the new list item.
         (if (eq direction 'above)
             (org-beginning-of-item)
           (end-of-line))
         (let* ((ctx-item? (eq 'item (org-element-type context)))
                (ctx-cb (org-element-property :contents-begin context))
                ;; Hack to handle edge case where the point is at the
                ;; beginning of the first item
                (beginning-of-list? (and (not ctx-item?)
                                         (= ctx-cb orig-point)))
                (item-context (if beginning-of-list?
                                  (org-element-context)
                                context))
                ;; Horrible hack to handle edge case where the
                ;; line of the bullet is empty
                (ictx-cb (org-element-property :contents-begin item-context))
                (empty? (and (eq direction 'below)
                             ;; in case contents-begin is nil, or contents-begin
                             ;; equals the position end of the line, the item is
                             ;; empty
                             (or (not ictx-cb)
                                 (= ictx-cb
                                    (1+ (point))))))
                (pre-insert-point (point)))
           ;; Insert dummy content, so that `org-insert-item'
           ;; inserts content below this item
           (when empty?
             (insert " "))
           (org-insert-item (org-element-property :checkbox context))
           ;; Remove dummy content
           (when empty?
             (delete-region pre-insert-point (1+ pre-insert-point))))))
      ;; Add a new table row
      ((or `table `table-row)
       (pcase direction
         ('below (save-excursion (org-table-insert-row t))
                 (org-table-next-row))
         ('above (save-excursion (org-shiftmetadown))
                 (+org/table-previous-row))))

      ;; Otherwise, add a new heading, carrying over any todo state, if
      ;; necessary.
      (_
       (let ((level (or (org-current-level) 1)))
         ;; I intentionally avoid `org-insert-heading' and the like because they
         ;; impose unpredictable whitespace rules depending on the cursor
         ;; position. It's simpler to express this command's responsibility at a
         ;; lower level than work around all the quirks in org's API.
         (pcase direction
           (`below
            (let (org-insert-heading-respect-content)
              (goto-char (line-end-position))
              (org-end-of-subtree)
              (insert "\n" (make-string level ?*) " ")))
           (`above
            (org-back-to-heading)
            (insert (make-string level ?*) " ")
            (save-excursion (insert "\n"))))
         (run-hooks 'org-insert-heading-hook)
         (when-let* ((todo-keyword (org-element-property :todo-keyword context))
                     (todo-type    (org-element-property :todo-type context)))
           (org-todo
            (cond ((eq todo-type 'done)
                   ;; Doesn't make sense to create more "DONE" headings
                   (car (+org-get-todo-keywords-for todo-keyword)))
                  (todo-keyword)
                  ('todo)))))))

    (when (org-invisible-p)
      (org-show-hidden-entry))
    (when (and (bound-and-true-p evil-local-mode)
               (not (evil-emacs-state-p)))
      (evil-insert 1))))

;;;###autoload
(defun +org/insert-item-below (count)
  "Inserts a new heading, table cell or item below the current one."
  (interactive "p")
  (dotimes (_ count) (+org--insert-item 'below)))

;;;###autoload
(defun +org/insert-item-above (count)
  "Inserts a new heading, table cell or item above the current one."
  (interactive "p")
  (dotimes (_ count) (+org--insert-item 'above)))


(defun org-make-olist (arg)
  (interactive "P")
  (let ((n (or arg 1)))
    (when (region-active-p)
      (setq n (count-lines (region-beginning)
                           (region-end)))
      (goto-char (region-beginning)))
    (dotimes (i n)
      (beginning-of-line)
      (insert (concat (number-to-string (1+ i)) ". "))
      (forward-line))))
\end{minted}
\section{Highlight yanked text}
\label{sec:org0a2e5c3}
This function describes itself so\ldots{}
\begin{minted}[linenos=true,numbersep=2pt,breaklines=true,frame=leftline,framerule=2pt,labelposition=bottomline,bgcolor=jpyellow!70]{common-lisp}
(defun meain/evil-yank-advice (orig-fn beg end &rest args)
  (pulse-momentary-highlight-region beg end 'mode-line-active)
  (apply orig-fn beg end args))
(advice-add 'evil-yank :around 'meain/evil-yank-advice)
\end{minted}
\section{Replace '-' with '•'.}
\label{sec:org59e525d}
\begin{minted}[linenos=true,numbersep=2pt,breaklines=true,frame=leftline,framerule=2pt,labelposition=bottomline,bgcolor=jpyellow!70]{common-lisp}
(defun jp/org-font-setup ()
  ;; Replace list hyphen with dot
  (font-lock-add-keywords 'org-mode
                          '(("^ *\\([-]\\) "
                             (0 (prog1 () (compose-region (match-beginning 1) (match-end 1) "•")))))))

(jp/org-font-setup)
\end{minted}
\section{Custom\textsubscript{id}.}
\label{sec:org15e0341}
\begin{minted}[linenos=true,numbersep=2pt,breaklines=true,frame=leftline,framerule=2pt,labelposition=bottomline,bgcolor=jpyellow!70]{elisp}
;;;; org-id
(declare-function org-id-add-location "org")
(declare-function org-with-point-at "org")
(declare-function org-entry-get "org")
(declare-function org-id-new "org")
(declare-function org-entry-put "org")

;; Copied from this article (with minor tweaks from my side):
;; https://writequit.org/articles/emacs-org-mode-generate-ids.html.
(defun jp/org--id-get (&optional pom create prefix)
  "Get the CUSTOM_ID property of the entry at point-or-marker POM.
If POM is nil, refer to the entry at point.  If the entry does
not have an CUSTOM_ID, the function returns nil.  However, when
CREATE is non nil, create a CUSTOM_ID if none is present already.
PREFIX will be passed through to `org-id-new'.  In any case, the
CUSTOM_ID of the entry is returned."
  (org-with-point-at pom
    (let ((id (org-entry-get nil "CUSTOM_ID")))
      (cond
       ((and id (stringp id) (string-match \\S- id))
        id)
       (create
        (setq id (org-id-new (concat prefix "h")))
        (org-entry-put pom "CUSTOM_ID" id)
        (org-id-add-location id (format "%s" (buffer-file-name (buffer-base-buffer))))
        id)))))

(declare-function org-map-entries "org")

;;;###autoload
(defun jp/org-id-headlines ()
  "Add missing CUSTOM_ID to all headlines in current file."
  (interactive)
  (org-map-entries
   (lambda () (jp/org--id-get (point) t))))

;;;###autoload
(defun jp/org-id-headline ()
  "Add missing CUSTOM_ID to headline at point."
  (interactive)
  (jp/org--id-get (point) t))
\end{minted}
\section{Add ID's to headers - Org roam.}
\label{sec:orgf2fa6d4}
\begin{minted}[linenos=true,numbersep=2pt,breaklines=true,frame=leftline,framerule=2pt,labelposition=bottomline,bgcolor=jpyellow!70]{common-lisp}
(defun jp/org-id-store-link-for-headers ()
  "Run `org-id-store-link' for each header in the current buffer."
  (interactive)
  (save-excursion
    (goto-char (point-min))
    (while (re-search-forward org-heading-regexp nil t)
      (org-id-store-link))))
\end{minted}
\section{Toggle org-emphasis-markers}
\label{sec:orgdc8dec8}
\begin{minted}[linenos=true,numbersep=2pt,breaklines=true,frame=leftline,framerule=2pt,labelposition=bottomline,bgcolor=jpyellow!70]{elisp}
(defun jp/org-toggle-emphasis-markers (&optional arg)
  "Toggle emphasis markers and display a message."
  (interactive "p")
  (let ((markers org-hide-emphasis-markers)
        (msg ""))
    (when markers
      (setq-local org-hide-emphasis-markers nil)
      (setq msg "Emphasis markers are now visible."))
    (unless markers
      (setq-local org-hide-emphasis-markers t)
      (setq msg "Emphasis markers are now hidden."))
    (message "%s" msg)
    (when arg
      (font-lock-fontify-buffer))))
\end{minted}
\section{Insert Org's files to Outlook.}
\label{sec:org8bf0241}
\subsection{Export file to clipboard.}
\label{sec:orgc140211}
\begin{minted}[linenos=true,numbersep=2pt,breaklines=true,frame=leftline,framerule=2pt,labelposition=bottomline,bgcolor=jpyellow!70]{elisp}
(defun export-org-email ()
  "Export the current email org buffer and copy it to the
clipboard"
  (interactive)
  (let ((org-export-show-temporary-export-buffer nil)
        (org-html-head (org-email-html-head)))
    (org-html-export-as-html)
    (with-current-buffer "*Org HTML Export*"
      (kill-new (buffer-string)))
    (message "HTML copied to clipboard")))
(global-set-key (kbd "C-c C-x C-e") 'export-org-email)
\end{minted}
\subsection{Add some CSS to the file.}
\label{sec:org5ecc052}
\begin{minted}[linenos=true,numbersep=2pt,breaklines=true,frame=leftline,framerule=2pt,labelposition=bottomline,bgcolor=jpyellow!70]{elisp}
(defun org-email-html-head ()
  "Create the header with CSS for use with email"
  (concat
   "<style type=\"text/css\">\n"
   "<!--/*--><![CDATA[/*><!--*/\n"
   (with-temp-buffer
     (insert-file-contents
      "~/.emacs.d/src/css/org2outlook.css")
     (buffer-string))
   "/*]]>*/-->\n"
   "</style>\n"))
\end{minted}
\section{Org-failure-report.}
\label{sec:orgaa8e7bb}
\begin{minted}[linenos=true,numbersep=2pt,breaklines=true,frame=leftline,framerule=2pt,labelposition=bottomline,bgcolor=jpyellow!70]{common-lisp}
(defvar report-file "~/Desktop/report.org"
  "Path to the Org-mode file to store the failure/solution report.")

(defun report-failure-solution ()
  "Prompt for machine number, failure, solution, and repair time. Append to an Org-mode report file."
  (interactive)
  (let ((continue-loop t))
    (catch 'exit
      (while continue-loop
        (let* ((machine-number (read-string "Enter machine number (or 'exit' to finish): "))
               (lowercase-machine (downcase machine-number)))
          (when (equal lowercase-machine "exit")
            (message "Exiting failure/solution report.")
            (setq continue-loop nil)
            (throw 'exit nil))

          (let* ((failure (read-string "Enter failure: "))
                 (solution (read-string "Enter solution: "))
                 (repair-time (read-string "Enter repair time (e.g., 2 hours): "))
                 (marker (concat "Machine " machine-number)))

            (with-temp-buffer
              (insert-file-contents report-file)
              (goto-char (point-min))

              (if (re-search-forward marker nil t)
                  (progn
                    (goto-char (line-end-position))
                    (insert "\n*** Failure: " failure "\n")
                    (insert "    - Solution: " solution "\n")
                    (insert "    - Repair Time: " repair-time "\n"))
                (goto-char (point-max))
                (insert "\n* " marker "\n")
                (insert "** Failure: " failure "\n")
                (insert "   - Solution: " solution "\n")
                (insert "   - Repair Time: " repair-time "\n"))

              (write-region (point-min) (point-max) report-file nil 'append))))))))

(global-set-key (kbd "C-c <f12>") 'report-failure-solution)
\end{minted}
\section{Org TODO auto update.}
\label{sec:orga118731}
\begin{minted}[linenos=true,numbersep=2pt,breaklines=true,frame=leftline,framerule=2pt,labelposition=bottomline,bgcolor=jpyellow!70]{common-lisp}
(defun org-todo-if-needed (state)
  "Change header state to STATE unless the current item is in STATE already."
  (unless (string-equal (org-get-todo-state) state)
    (org-todo state)))

(defun ct/org-summary-todo-cookie (n-done n-not-done)
  "Switch header state to DONE when all subentries are DONE, to TODO when none are DONE, and to DOING otherwise"
  (let (org-log-done org-log-states)   ; turn off logging
    (org-todo-if-needed (cond ((= n-done 0)
                               "TODO")
                              ((= n-not-done 0)
                               "DONE")
                              (t
                               "DOING")))))
(add-hook 'org-after-todo-statistics-hook #'ct/org-summary-todo-cookie)

(defun ct/org-summary-checkbox-cookie ()
  "Switch header state to DONE when all checkboxes are ticked, to TODO when none are ticked, and to DOING otherwise"
  (let (beg end)
    (unless (not (org-get-todo-state))
      (save-excursion
        (org-back-to-heading t)
        (setq beg (point))
        (end-of-line)
        (setq end (point))
        (goto-char beg)
        ;; Regex group 1: %-based cookie
        ;; Regex group 2 and 3: x/y cookie
        (if (re-search-forward "\\[\\([0-9]*%\\)\\]\\|\\[\\([0-9]*\\)/\\([0-9]*\\)\\]"
                               end t)
            (if (match-end 1)
                ;; [xx%] cookie support
                (cond ((equal (match-string 1) "100%")
                       (org-todo-if-needed "DONE"))
                      ((equal (match-string 1) "0%")
                       (org-todo-if-needed "TODO"))
                      (t
                       (org-todo-if-needed "DOING")))
              ;; [x/y] cookie support
              (if (> (match-end 2) (match-beginning 2)) ; = if not empty
                  (cond ((equal (match-string 2) (match-string 3))
                         (org-todo-if-needed "DONE"))
                        ((or (equal (string-trim (match-string 2)) "")
                             (equal (match-string 2) "0"))
                         (org-todo-if-needed "TODO"))
                        (t
                         (org-todo-if-needed "DOING")))
                (org-todo-if-needed "DOING"))))))))
(add-hook 'org-checkbox-statistics-hook #'ct/org-summary-checkbox-cookie)
\end{minted}
\section{Org-buffer-scratchpad.}
\label{sec:org415cfb8}
\begin{minted}[linenos=true,numbersep=2pt,breaklines=true,frame=leftline,framerule=2pt,labelposition=bottomline,bgcolor=jpyellow!70]{common-lisp}
(defun new-scratch-pad ()
  "Create a new org-mode buffer for random stuff."
  (interactive)
  (progn
    (let ((buffer (generate-new-buffer "Org-scratch-buffer")))
      (switch-to-buffer buffer)
      (setq buffer-offer-save t)
      (org-mode)
      (olivetti-mode t))))
\end{minted}
\section{Toggle buffers.}
\label{sec:org70ce57a}
\begin{minted}[linenos=true,numbersep=2pt,breaklines=true,frame=leftline,framerule=2pt,labelposition=bottomline,bgcolor=jpyellow!70]{common-lisp}
;; Toggle *scratch* buffer.
(defun toggle-scratch-buffer ()
  "Toggle the *scratch* buffer"
  (interactive)
  (if (string= (buffer-name) "*scratch*")
      (bury-buffer)
    (switch-to-buffer (get-buffer-create "*scratch*"))))

(defun toggle-org-buffer ()
  "Toggle the Org-scratch-buffer buffer"
  (interactive)
  (if (equal (buffer-name (current-buffer)) "Org-scratch-buffer")
      (if (one-window-p t)
          (switch-to-buffer (other-buffer))
        (delete-window))
    (if (get-buffer "Org-scratch-buffer")
        (if (get-buffer-window "Org-scratch-buffer")
            (progn
              (bury-buffer "Org-scratch-buffer")
              (delete-window (get-buffer-window "Org-scratch-buffer")))
          (switch-to-buffer "Org-scratch-buffer"))
      (new-scratch-pad))))

;; Toggle *eshell* buffer.
(defun toggle-eshell-buffer ()
  "Toggle the *eshell* buffer"
  (interactive)
  (if (string= (buffer-name) "*eshell*")
      (bury-buffer)
    (switch-to-buffer (get-buffer-create "*eshell*"))))
\end{minted}
\section{Doc at point.}
\label{sec:org093de8d}
\begin{minted}[linenos=true,numbersep=2pt,breaklines=true,frame=leftline,framerule=2pt,labelposition=bottomline,bgcolor=jpyellow!70]{common-lisp}
(defun my-show-doc-or-describe-symbol ()
  "Show LSP UI doc if LSP is active, otherwise describe symbol at point."
  (interactive)
  (if (bound-and-true-p lsp-mode)
      (lsp-ui-doc-glance)
    (describe-symbol-at-point)))
\end{minted}
\section{Duplicate \& move up/down lines.}
\label{sec:org260b5d6}
\begin{minted}[linenos=true,numbersep=2pt,breaklines=true,frame=leftline,framerule=2pt,labelposition=bottomline,bgcolor=jpyellow!70]{common-lisp}
(defun duplicate-line ()
  (interactive)
  (let ((line-text (thing-at-point 'line t)))
    (save-excursion
      (move-end-of-line 1)
      (newline)
      (insert line-text)))
  (forward-line 1))

(defun move-line-up ()
  (interactive)
  (when (not (= (line-number-at-pos) 1))
    (transpose-lines 1)
    (forward-line -2)))

(defun move-line-down ()
  (interactive)
  (forward-line 1)
  (when (not (= (line-number-at-pos) (point-max)))
    (transpose-lines 1))
  (forward-line -1))
\end{minted}
\section{Hide passwords in org files.}
\label{sec:orga0e5be5}
\begin{minted}[linenos=true,numbersep=2pt,breaklines=true,frame=leftline,framerule=2pt,labelposition=bottomline,bgcolor=jpyellow!70]{common-lisp}
;; Define a custom face for the highlight
(defface my-highlight-face
  '((t (:foreground "gray"))) ; Change "red" to your desired color
  "Face for highlighting !!word!! patterns.")

;; Function to replace matched text with asterisks
(defun replace-with-asterisks (limit)
  "Replace !!word!! with asterisks up to LIMIT."
  (while (re-search-forward "!!\\(.*?\\)!!" limit t)
    (let* ((match (match-string 1))
           (start (match-beginning 0))
           (end (match-end 0))
           (asterisks (make-string (length match) ?*)))
      (add-text-properties start end `(display ,asterisks face my-highlight-face)))))

;; Add custom keyword for font-lock in org-mode
(defun my/org-mode-custom-font-lock ()
  "Add custom font-lock keywords for org-mode."
  (font-lock-add-keywords nil
                          '((replace-with-asterisks))))

;; Hook the custom font-lock configuration into org-mode
(add-hook 'org-mode-hook 'my/org-mode-custom-font-lock)
\end{minted}
\section{Help at point.}
\label{sec:orgda0ff03}
\begin{minted}[linenos=true,numbersep=2pt,breaklines=true,frame=leftline,framerule=2pt,labelposition=bottomline,bgcolor=jpyellow!70]{common-lisp}
(defun describe-symbol-at-point ()
  "Display the documentation of the symbol at point, if it exists."
  (interactive)
  (let ((symbol (symbol-at-point)))
    (if symbol
        (cond
         ((fboundp symbol) (describe-function symbol))
         ((boundp symbol) (describe-variable symbol))
         (t (message "No documentation found for symbol at point: %s" symbol)))
      (message "No symbol at point"))))
\end{minted}
\section{Custom functions.}
\label{sec:orge26144a}
\begin{minted}[linenos=true,numbersep=2pt,breaklines=true,frame=leftline,framerule=2pt,labelposition=bottomline,bgcolor=jpyellow!70]{common-lisp}
;; Open files in the lisp folder
(require 'find-lisp)
(defun open-lisp-and-org-files ()
  "Open a Lisp or Org file from ~/.emacs.d/lisp directory, including subfolders."
  (interactive)
  (let* ((directory "~/.emacs.d/lisp")
         (el-files (find-lisp-find-files directory ".*\\.el$"))
         (org-files (find-lisp-find-files directory ".*\\.org$"))
         (all-files (append el-files org-files))
         (file (completing-read "Select file: " all-files nil t)))
    (find-file file)))

(defun open-post-org-files ()
  "List and open Org files in the ~/repos/jpacheco.xyz/posts/ directory."
  (interactive)
  (let* ((org-directory "/home/javier/repos/jpacheco.xyz/content/posts/")
         (org-files (directory-files org-directory nil "\\.org$"))
         (chosen-file (completing-read "Choose an Org file: " org-files nil t)))
    (find-file (expand-file-name chosen-file org-directory))))

;; Follow urls in the buffer
(defun list-and-open-url-in-buffer ()
  "List all URLs in the current buffer, display them in the minibuffer, and open a selected URL in the browser."
  (interactive)
  (let (urls)
    (save-excursion
      (goto-char (point-min))
      (while (re-search-forward "\\(http\\|https\\|ftp\\|file\\|mailto\\):[^ \t\n]+" nil t)
        (push (match-string 0) urls)))
    (if urls
        (let ((url (completing-read "Select URL to open: " (reverse urls) nil t)))
          (browse-url url))
      (message "No URLs found in the buffer."))))

;; Export org files to pdf using tectonic
(defun org-export-to-latex-and-compile-with-tectonic ()
  "Export the current Org file to LaTeX, then compile with tectonic using shell-escape."
  (interactive)
  (let* ((org-file (buffer-file-name))
         (tex-file (concat (file-name-sans-extension org-file) ".tex"))
         (tectonic-command (concat "tectonic -Z shell-escape " tex-file)))
    ;; Export Org file to LaTeX
    (org-latex-export-to-latex)
    ;; Run tectonic command in a temporary buffer to avoid displaying the output
    (with-temp-buffer
      (shell-command tectonic-command (current-buffer)))
    (message "Compiled %s to PDF with Tectonic." tex-file)))

(global-set-key (kbd "C-c e l") 'org-export-to-latex-and-compile-with-tectonic)

;; Update my web-page
(defun run-build-script ()
  "Run build.sh script to update my web page"
  (interactive)
  (let ((emacs-path (executable-find "emacs"))
        (with-temp-buffer
          (cd "~/repos/jpacheco.xyz"))
        (script-path "~/repos/jpacheco.xyz/build.sh"))
    (message "Starting the process to run build-site.el...")
    (if (file-exists-p script-path)
        (progn
          (message "Emacs executable found at: %s" emacs-path)
          (message "Script found at: %s" script-path)
          (async-shell-command "sh build.sh"))
      (message "Script not found: %s" script-path))))
(global-set-key (kbd "C-c u p") 'run-build-script)
\end{minted}
\section{Provide this code:}
\label{sec:orgaa27c14}
\begin{minted}[linenos=true,numbersep=2pt,breaklines=true,frame=leftline,framerule=2pt,labelposition=bottomline,bgcolor=jpyellow!70]{common-lisp}
(provide 'utilities)
\end{minted}
\end{document}
